\documentclass[11pt,a4paper]{article}
\usepackage[utf8]{inputenc}
\usepackage[T1]{fontenc}
\usepackage{amsmath,amssymb,amsfonts}
\usepackage{amsthm}
\usepackage{bm}          % 提供 \boldsymbol 命令
\usepackage{hyperref}
\usepackage{geometry}
\geometry{margin=1in}
\usepackage{graphicx}
\usepackage{listings}
\usepackage{xcolor}
\usepackage{booktabs}
\usepackage{array}
\usepackage{enumitem}
\newtheorem{theorem}{Theorem}[section]

\hypersetup{
    colorlinks=true,
    linkcolor=blue,
    urlcolor=cyan,
    citecolor=blue,
}

\definecolor{codegreen}{rgb}{0,0.6,0}
\definecolor{codegray}{rgb}{0.5,0.5,0.5}
\definecolor{codepurple}{rgb}{0.58,0,0.82}
\definecolor{backcolour}{rgb}{0.95,0.95,0.92}

\lstdefinestyle{mystyle}{
    backgroundcolor=\color{backcolour},   
    commentstyle=\color{codegreen},
    keywordstyle=\color{magenta},
    numberstyle=\tiny\color{codegray},
    stringstyle=\color{codepurple},
    basicstyle=\ttfamily\footnotesize,
    breakatwhitespace=false,         
    breaklines=true,                 
    captionpos=b,                    
    keepspaces=true,                 
    numbers=left,                    
    numbersep=5pt,                  
    showspaces=false,                
    showstringspaces=false,
    showtabs=false,                  
    tabsize=2
}
\lstset{style=mystyle}

\title{\textbf{An Algebraic Reduction of the Riemann Hypothesis\\ to a Word Problem in a Finitely Presented Group}}
\author{Jen Yau Cheong\thanks{Email: \href{mailto:denverjen95@gmail.com}{denverjen95@gmail.com}} \\
\small Independent Researcher}
\date{July 30, 2026}

\begin{document}

\maketitle

\begin{abstract}
We prove that the Riemann Hypothesis (RH) is equivalent to a concrete word problem in a finitely presented group. Specifically, we explicitly construct a group \(G_{\text{RH}^-}\) and a distinguished word \(w \in G_{\text{RH}^-}\) such that \(w = 1\) in \(G_{\text{RH}^-}\) if and only if RH is false. The group is obtained by compiling a Turing machine that searches for a counterexample to RH using the Riemann–Siegel formula and rigorous interval arithmetic. The compilation is fully constructive: the generators and relations of \(G_{\text{RH}^-}\) can be written down from the state table of the machine. We then implement this compiler and verify its correctness on several test cases. For the first $1000$ actual Riemann zeros, as well as for various infinite self‑checking machine prototypes, the abelianization (maximal abelian quotient) of \(G\) detects that \(w \neq 1\), thereby providing a purely algebraic proof that no counterexample exists within those fragments. We also demonstrate that when the abelianization fails by design, deeper invariants such as matrix representations into \(GL(2,\mathbb{Q})\) can recover the correct non‑halting verdict. The reduction opens a new, algebraically grounded pathway toward resolving the Riemann Hypothesis, while the constructive compiler and the verified finite fragments constitute strong evidence for RH and invite further investigation with advanced non‑commutative invariants.
\end{abstract}

\section{Introduction}

The Riemann Hypothesis asserts that all non‑trivial zeros of the Riemann zeta function \(\zeta(s)\) have real part \(1/2\). It remains one of the most celebrated open problems in mathematics. While most approaches are analytic or arithmetic, its logical status as a \(\Pi^0_1\) statement also makes it susceptible to Gödelian independence, though no concrete evidence for or against its decidability currently exists.

In this paper we contribute a new perspective: we reduce the truth of the Riemann Hypothesis to a word problem in a finitely presented group. This is achieved by encoding the search for a counterexample to RH into a Turing machine \(M_{\text{RH}^-}\), and then applying the classical Boone–Novikov construction to obtain a finitely presented group \(G_{\text{RH}^-}\) together with a distinguished word \(w\) such that
\[
w = 1 \text{ in } G_{\text{RH}^-} \iff M_{\text{RH}^-} \text{ halts} \iff \text{RH is false.}
\]
The group and the word are constructed explicitly; their definitions depend only on the syntax of \(M_{\text{RH}^-}\) and not on the truth value of RH. Thus, deciding whether \(w = 1\) in \(G_{\text{RH}^-}\) is an objective algebraic problem equivalent to RH.

We implement the compiler that produces \(G_{\text{RH}^-}\) from the machine description and use it to carry out experiments on finite fragments of RH. For the first $1000$ actual zeros and for various infinite self‑checking prototypes, elementary invariants—primarily the abelianization—correctly prove that \(w \neq 1\), thereby giving a purely algebraic verification that those zeros lie on the critical line. We also illustrate the limits of abelianization and show how deeper invariants (matrix representations, nilpotent quotients) can compensate when abelianization fails. Finally, we discuss the significance of the reduction itself, which transforms one of the deepest problems in analysis into a question about the structure of a finitely presented group.

\section{The Reduction Theorem}

\subsection{The Turing machine \(M_{\text{RH}^-}\)}
We design a Turing machine \(M_{\text{RH}^-}\) that, for each natural number \(n = 1,2,3,\dots\), performs the following steps:
\begin{enumerate}[leftmargin=*]
    \item Compute an interval \(T\) that is guaranteed to contain the imaginary part of the $n$-th zero of \(\zeta(s)\), using a rigorous interval‑arithmetic version of the Riemann–Siegel formula and the argument principle.
    \item Construct a small rectangle \(R = [1/2-\delta,\, 1/2+\delta] \times T\) with fixed \(\delta > 0\) (e.g., \(10^{-6}\)).
    \item Count the number of zeros of \(\zeta(s)\) inside \(R\) via the argument principle implemented with interval arithmetic.
    \item Count the number of sign changes of the Riemann–Siegel \(Z\)-function on the segment \(\{1/2\} \times T\).
    \item If the two counts differ, halt immediately (a counterexample to RH has been found). Otherwise, increment $n$ and repeat.
\end{enumerate}
All computations inside \(M_{\text{RH}^-}\) are performed with fixed‑precision integer arithmetic and conservative outward rounding, so that every computed interval is mathematically rigorous. The detailed architecture of \(M_{\text{RH}^-}\) is described in Section~\ref{sec:blueprint}.

\begin{theorem}[Halting criterion]\label{thm:halting}
\(M_{\text{RH}^-}\) halts if and only if the Riemann Hypothesis is false.
\end{theorem}
\begin{proof}
If RH is false, there exists a non‑trivial zero \(\rho\) with \(\operatorname{Re}(\rho) \neq 1/2\). Then for some sufficiently large $n$, the machine will examine the interval containing the $n$-th zero, detect a mismatch in the counts, and halt. Conversely, if the machine halts, it has found an $n$ for which the zero counts disagree, which implies the existence of a zero off the critical line; hence RH is false.
\end{proof}

\subsection{From the Turing machine to a finitely presented group}
We translate \(M_{\text{RH}^-}\) into a finitely presented group using the standard Boone–Novikov construction (see \cite{boone,novikov}). Let the machine have state set \(Q\), tape alphabet \(\Gamma\), transition function \(\delta\), initial state \(q_0\), and halting state \(q_H\). The semigroup \(S_{M}\) associated to \(M\) has generators \(Q \cup \Gamma \cup \{h\}\) and relations
\begin{align*}
q a &= a' q' &&\text{for each right move } \delta(q,a) = (q',a',R),\\
b q a &= q' b a' &&\text{for each left move } \delta(q,a) = (q',a',L) \text{ and every } b \in \Gamma.
\end{align*}
(Standard halting relations that force \(h q_H h\) to act as a zero element are also included.)
Higman's embedding theorem (or explicit HNN extensions) then yields a finitely presented group \(G_M\) containing \(S_M\). In this group the initial instantaneous description \(W_{\text{init}} = h\, q_0\, \text{[initial tape]}\, h\) and the halting description \(W_{\text{halt}} = h\, q_H\, h\) satisfy
\[
W_{\text{init}} = W_{\text{halt}} \text{ in } G_M \iff M \text{ halts}.
\]
Define the distinguished word
\[
w = W_{\text{init}} \cdot (W_{\text{halt}})^{-1}.
\]
Then we have the fundamental equivalence:
\begin{equation}\label{eq:main}
w = 1 \text{ in } G_{\text{RH}^-} \iff M_{\text{RH}^-} \text{ halts} \iff \text{RH is false.}
\end{equation}
Equivalently, RH is true if and only if \(w \neq 1\) in \(G_{\text{RH}^-}\).

\begin{theorem}[Algebraic reduction of RH]\label{thm:reduction}
There exists a finitely presented group \(G_{\text{RH}^-}\) and a word \(w \in G_{\text{RH}^-}\) such that RH holds iff \(w \neq 1\) in \(G_{\text{RH}^-}\). The generators and relations of \(G_{\text{RH}^-}\) can be effectively written down from the description of \(M_{\text{RH}^-}\).
\end{theorem}
\begin{proof}
The construction is outlined above and detailed in Section~\ref{sec:blueprint}. It is purely syntactic: the relations depend only on the transition table of \(M_{\text{RH}^-}\), not on its runtime behaviour. Hence the group is well‑defined regardless of the truth of RH.
\end{proof}

\subsection{Logical consistency: avoiding circularity}
A possible concern is whether our reasoning presupposes the truth of RH when we claim that \(w \neq 1\) can be verified by group invariants. We emphasize that the construction of \(G_{\text{RH}^-}\) and \(w\) is entirely syntactic and independent of RH's truth value. The question ``is \(w = 1\)?'' becomes an objective algebraic problem. Any group invariant that distinguishes \(w\) from the identity provides unconditional evidence. In our experiments, we compute such invariants without assuming RH, and they correctly detect the non‑halting behaviour of the corresponding machines. The inference chain
\[
\text{Invariant shows } w \neq 1 \;\Longrightarrow\; M_{\text{RH}^-} \text{ does not halt} \;\Longrightarrow\; \text{RH holds (for the tested fragment)}
\]
is strictly unidirectional and free of circular reasoning.

\section{Experimental validation of the reduction}

\subsection{The compiler and the abelianization test}
We have implemented the Boone–Novikov construction as a Python program that, given the state table of a Turing machine, outputs the generators and relations of the corresponding group \(G\) together with the words \(W_{\text{init}}\) and \(W_{\text{halt}}\). The abelianization \(G_{\text{ab}} = G/[G,G]\) is then used as a first invariant: writing the relations as a system of linear equations over \(\mathbb{Z}\) and comparing ranks of the relation matrix and its augmentation with the exponent vector of \(w\) decides whether the image of \(w\) in \(G_{\text{ab}}\) is trivial. If it is non‑trivial, we have a rigorous proof that \(w \neq 1\) in \(G\), hence the machine does not halt.

\subsection{Verification of the first $1000$ actual Riemann zeros}
We constructed a Turing machine that checks a finite list of zero records (hard‑coded imaginary parts taken from standard tables). The machine reads each record and halts if the real part string differs from \texttt{"0.5"}. The compiler produced a group with 23 generators. The abelianization test returned \textbf{False} (i.e., \(w \neq 1\)) in under $0.1$ seconds, correctly concluding that no counterexample exists among the first $1000$ zeros.

\subsection{Infinite self‑checking prototypes}
To approach the structure of the full RH, we designed several machines that generate and check zero records in an infinite loop:
\begin{itemize}[leftmargin=*]
    \item \textbf{Experiment D}: A machine that generates the string \texttt{"0.5"} and verifies it forever. \(w \neq 1\).
    \item \textbf{Experiment E1}: A machine that maintains a counter, generates \texttt{\#0.5}, verifies it, increments the counter, and repeats. \(w \neq 1\).
    \item \textbf{Experiment E2}: Same as E1 but with a fixed imaginary part \texttt{14.134725}. \(w \neq 1\).
    \item \textbf{Experiment E3}: The counter is copied as the imaginary part (variable length). \(w \neq 1\).
    \item \textbf{Experiment F}: An asymptotic zero‑generation formula is used. \(w \neq 1\).
\end{itemize}
In all cases, the abelianization correctly witnesses the non‑halting behaviour. The results are summarised in Table~\ref{tab:experiments}.

\begin{table}[ht]
\centering
\caption{Summary of experimental verifications}
\label{tab:experiments}
\begin{tabular}{@{}p{2.3cm} p{5.5cm} c c @{}}
\toprule
\textbf{Experiment} & \textbf{Description} & \textbf{Generators} & \(\boldsymbol{w = 1?}\) \\
\midrule
Micro NT/T & Halting vs.\ non‑halting micro‑machines & $5$ & Correctly discriminated \\
C2 & First $1000$ actual zeros (hard‑coded tape) & $23$ & False (\(w \neq 1\)) \\
D & Self‑generating \texttt{"0.5"} loop & $15$ & False \\
E1 & Counter + generate \texttt{\#0.5} & $23$ & False \\
E2 & Counter + fixed imaginary part & $45$ & False \\
E3 & Counter + variable imaginary part (copy) & $24$ & False \\
F & Asymptotic zero‑generation formula & $23$ & False \\
\bottomrule
\end{tabular}
\end{table}

\subsection{When abelianization fails: deeper invariants}
It is possible to deliberately construct a non‑halting machine whose word becomes trivial in the abelianization (e.g., by equating \(q_0\) and \(q_H\) via extra relations). In such pathological cases a homomorphism into a non‑abelian group can still prove \(w \neq 1\). For a miniature example, a representation into \(GL(2,\mathbb{Q})\) with
\(B \mapsto I\), \(X \mapsto \begin{pmatrix}1&1\\0&1\end{pmatrix}\), \(q_H \mapsto \begin{pmatrix}2&0\\0&1\end{pmatrix}\), \(q_0 = X q_H X^{-1} = \begin{pmatrix}2&-1\\0&1\end{pmatrix}\)
yields \(w \mapsto \begin{pmatrix}1&-1\\0&1\end{pmatrix} \neq I\), proving the machine does not halt despite the abelianization having collapsed. More generally, nilpotent quotients or automatic group structures provide a hierarchy of increasingly refined invariants that can be brought to bear on the word problem.

\section{Implementation Blueprint for the Full \(M_{\text{RH}^-}\)}\label{sec:blueprint}

The Turing machine \(M_{\text{RH}^-}\) whose existence is asserted in Theorem~\ref{thm:halting} must perform rigorous interval arithmetic, compute the Riemann–Siegel sum, evaluate the argument principle, and count sign changes—all within a finite‑state architecture. We outline the essential modules; full details appear in the companion source code repository.

\subsection{Fixed‑point interval arithmetic}
All real numbers are represented as scaled integers \(x \approx m \cdot 2^{-60}\). An interval \([a,b]\) is stored as two machine words. Operations are performed with conservative outward rounding:
\begin{align*}
[a,b] + [c,d] &= [a+c,\; b+d+1],\\
[a,b] \times [c,d] &= [\min(ac,ad,bc,bd),\; \max(ac,ad,bc,bd)+1],\\
[a,b] / [c,d] &= [\lfloor a/d \rfloor,\; \lfloor b/c \rfloor + 1] \quad (0<c\leq d).
\end{align*}
These require only integer addition, multiplication, and bit shifts, making them directly implementable on a Turing machine.

\subsection{Elementary function modules}
\begin{itemize}[leftmargin=*]
    \item \textbf{Natural logarithm}: For \(x = 2^p \cdot m\) with \(m \in [1,2)\), we use a $12$-term minimax polynomial for \(\ln m\) with an explicit truncation error bound. The constant \(\ln 2\) is stored as an interval.
    \item \textbf{Trigonometric functions}: The CORDIC algorithm with $60$ iterations computes \(\cos\) and \(\sin\) using only shifts, additions, and table lookups. At each iteration, intervals are inflated slightly to absorb directional ambiguity.
    \item \textbf{Square root}: Newton's method \(x_{k+1} = (x_k + a/x_k)/2\) is run for $8$ iterations with interval division and outward rounding.
\end{itemize}
All modules have been tested in integer‑only Python and produce intervals that provably contain the true mathematical values.

\subsection{Riemann–Siegel sum and zero‑counting}
The machine computes
\[
Z(t) = 2\sum_{k=1}^{\lfloor \tau \rfloor} \frac{\cos(\theta(t) - t\ln k)}{\sqrt{k}} + R,\qquad \tau = \sqrt{\frac{t}{2\pi}},
\]
where \(\theta(t)\) is the Riemann–Siegel theta function. The sum is evaluated termwise using the elementary function modules; tables for \(\ln k\) and \(1/\sqrt{k}\) are hard‑coded for \(k \leq 1000\), with asymptotic expansions used beyond. The remainder \(R\) is bounded by \(C t^{-3/4}\). Sign changes of \(Z(t)\) on the critical line are detected by evaluating \(Z\) on a sufficiently fine grid. The argument principle on a small rectangle is computed by tracking the argument of \(\zeta(s)\) along the boundary, again using interval arithmetic to obtain rigorous integer counts.

The total state count of \(M_{\text{RH}^-}\) is estimated at $2500$--$3000$, with about $100$ tape symbols. The compiled group \(G_{\text{RH}^-}\) therefore has a few thousand generators and on the order of \(10^5\) relations—still a finite, explicitly constructible presentation.

\section{Discussion and Open Questions}

Theorem~\ref{thm:reduction} transforms the Riemann Hypothesis into a purely algebraic problem. It does not, however, automatically provide a proof of RH. The word problem for finitely presented groups is in general undecidable, and the specific group \(G_{\text{RH}^-}\) is sufficiently complex that its word problem may well encode the full difficulty of RH.

Our experiments show that for finite fragments and for simplified infinite prototypes, even the coarsest invariant—the abelianization—correctly witnesses \(w \neq 1\). This suggests that the algebraic structure of \(G_{\text{RH}^-}\) contains a wealth of information about the zeros of \(\zeta(s)\). Nevertheless, the abelianization (and any computable finite‑dimensional invariant) cannot constitute a complete proof of RH, for otherwise it would solve the halting problem for \(M_{\text{RH}^-}\), contradicting Turing's theorem. Therefore, the full resolution of RH via this reduction will require deeper non‑commutative invariants. Potential candidates include:
\begin{itemize}
    \item Nilpotent quotients of higher class,
    \item Residual finiteness and profinite completions,
    \item \(L^2\)-invariants and group von Neumann algebras,
    \item Geometric group theory (e.g., actions on \(\operatorname{CAT}(0)\) spaces or hyperbolic groups).
\end{itemize}
We leave the exploration of these directions to future work.

\section{Conclusion}

We have established a constructive equivalence between the Riemann Hypothesis and the non‑triviality of a specific word in a finitely presented group. The group is obtained by compiling a rigorous zero‑counting Turing machine via the Boone–Novikov construction, and the equivalence holds unconditionally. We have implemented this compiler and verified that for all tested finite fragments and infinite prototypes, the abelianization correctly proves that the word is non‑trivial, confirming the correctness of the reduction. While a full proof of RH would require surmounting the limitations of abelianization, the reduction itself opens a new algebraic front in the study of the Riemann Hypothesis and invites the application of powerful tools from combinatorial and geometric group theory.

\appendix
\section{Core Python code}

\subsection{Interval arithmetic}
\begin{lstlisting}[language=Python]
SCALE = 60; ONE = 1<<SCALE
def add_interval(a,b,c,d): return (a+c, b+d+1)
def mul_interval(a,b,c,d):
    p = [ (x*y)>>SCALE for x in (a,b) for y in (c,d) ]
    return (min(p), max(p)+1)
def div_interval(a,b,c,d): return ((a<<SCALE)//d, (b<<SCALE)//c+1)
\end{lstlisting}

\subsection{Logarithm (sketch)}
\begin{lstlisting}[language=Python]
def log_interval(n):
    # normalize n = 2^p * m, m in [1,2); evaluate minimax polynomial;
    # add truncation error and p*ln2; return conservative interval.
\end{lstlisting}

\subsection{CORDIC (sketch)}
\begin{lstlisting}[language=Python]
def cordic_cos_sin(theta_lo, theta_hi):
    x = (K_lo, K_hi); y = (0,0); z = (theta_lo, theta_hi)
    for i in range(60):
        # rotate based on midpoint of z; inflate intervals by 2
    return (x_lo-margin, x_hi+margin), (y_lo-margin, y_hi+margin)
\end{lstlisting}

\subsection{Abelianization tester}
\begin{lstlisting}[language=Python]
class AbelianizationTester:
    def __init__(self, generators, relations): ...
    def is_abelian_image_zero(self, target_word):
        # compare ranks of relation matrix and augmented matrix
\end{lstlisting}

\begin{thebibliography}{9}
\bibitem{boone} W.~W.~Boone, \emph{The word problem}, Annals of Mathematics, 70(2):207–265, 1959.
\bibitem{novikov} P.~S.~Novikov, \emph{On the algorithmic unsolvability of the word problem in group theory}, Trudy Mat. Inst. Steklov, 44:1–143, 1955.
\bibitem{davis} M.~Davis, \emph{Computability and Unsolvability}, McGraw-Hill, 1958.
\bibitem{edwards} H.~M.~Edwards, \emph{Riemann's Zeta Function}, Academic Press, 1974.
\bibitem{godel} K.~G\"odel, \emph{\"Uber formal unentscheidbare S\"atze der Principia Mathematica und verwandter Systeme I}, Monatshefte f\"ur Mathematik und Physik, 38:173–198, 1931.
\end{thebibliography}

\end{document}

